% 本书原创：同一尺度绘出 beta=(1,0) 的精确水平流和两条数据曲线。
\begin{figure}[htbp]
\centering
\begin{tikzpicture}[x=1cm,y=1cm,>=stealth,every node/.style={font=\small},line width=.65pt]
\begin{scope}
\draw[->,gray!60] (-2.2,0)--(2.35,0) node[right] {\(x\)};
\draw[->,gray!60] (0,-1.7)--(0,1.85) node[above] {\(y\)};
\foreach \eta in {-1,0,1}{
  \draw[->,AnalysisGray] (-2.05,\eta)--(2.02,\eta);
}
\draw[->,AnalysisGray,very thick] (.4,1)--(1.4,1) node[midway,above=4pt] {\(\beta=(1,0)\)};
\draw[->,AnalysisBlue,very thick] (0,-1.5)--(0,1.5);
\node[AnalysisBlue,left] at (-.08,1.35) {\(\gamma\)};
\fill[AnalysisBlue] (0,0) circle(1.4pt);
\node at (0,2.6) {横截：\(\gamma(s)=(0,s)\)};
\node at (0,-2.15) {\(\det(\gamma',\beta)=-1\ne0\)};
\end{scope}
\begin{scope}[xshift=6.4cm]
\draw[->,gray!60] (-2.2,0)--(2.35,0) node[right] {\(x\)};
\draw[->,gray!60] (0,-1.7)--(0,1.85) node[above] {\(y\)};
\foreach \eta in {-1,0,1}{
  \draw[->,AnalysisGray] (-2.05,\eta)--(2.02,\eta);
}
\draw[->,AnalysisGray,very thick] (.4,1)--(1.4,1) node[midway,above=4pt] {\(\beta=(1,0)\)};
\draw[->,AnalysisBlue,very thick] (-1.5,0)--(1.5,0);
\node[AnalysisBlue,above] at (-.75,.08) {\(\gamma\)};
\fill[AnalysisBlue] (0,0) circle(1.4pt);
\node at (0,2.6) {相切：\(\gamma(s)=(s,0)\)};
\node at (0,-2.15) {\(\det(\gamma',\beta)=0\)};
\end{scope}
\end{tikzpicture}
\caption[数据曲线与特征轨道的横截和相切]{数据曲线与特征轨道的横截和相切。两图均取 \(\beta=(1,0)\)，灰线为水平轨道，粗灰箭头为长度为 \(1\) 的速度向量，蓝线为数据曲线；两图使用同一坐标尺度。左图的数据曲线与每条附近轨道唯一相交，\(\Psi(s,t)=(t,s)\) 可作局部坐标。右图的数据曲线完全落在一条轨道上，\(\Psi(s,t)=(s+t,0)\) 退化；它仅展示退化可能造成的失败，不代表任意行列式为零的情形都无解或不唯一。}
\label{b2:fig:07-noncharacteristic-data}
\end{figure}
